\chapter{外部数据快照的接入与核验}\label{app:shm}

\section{为什么单开一章}

2026 年 10 月，本项目获得了一份外部 A 股数据快照（四张 parquet，覆盖 2020-01 至 2026-08）。
它包含本地库完全没有的两样东西：\textbf{四档资金流}与\textbf{121 个财务指标}。

但"能拿到"不等于"能用"，更不等于"应该替换"。这一章记录四项核验的做法与结果，
以及两个更重要的结论：

\begin{enumerate}
  \item \textbf{看起来更全的数据，可能在一个更隐蔽的维度上更窄}——它缺的是最小那三成市值；
  \item \textbf{你以为需要它来修正的东西，可能你本来就已经很准}——本地估算的成交额
        中位误差只有 $0.3\%$。
\end{enumerate}

\section{数据结构}

\begin{center}
\begin{tabularx}{\textwidth}{@{}lrrX@{}}
\toprule
表 & 行数 & 列数 & 关键内容 \\
\midrule
\texttt{market\_daily} & 5,637,878 & 24 & 不复权与后复权 OHLC、成交额、VWAP/TWAP、换手率、流通股本 \\
\texttt{market\_moneyflow} & 5,637,848 & 33 & 超大单/大单/中单/小单的流入流出金额、量、笔数 \\
\texttt{valuation\_daily} & 9,553,741 & 15 & PE/PB/PS/PCF/PEG/股息率/资产负债率，日频 \\
\texttt{fundamental\_daily} & 5,637,878 & 123 & 121 个财务指标，含单季与 TTM \\
\bottomrule
\end{tabularx}
\end{center}

\section{三项核验}

\subsection{第一项：价格一致性}

把外部快照的收盘价与本地库（新浪通道，独立采集）按 \texttt{(code, date)} 内连接，
共匹配 \textbf{5,544,219} 行。

\begin{keybox}[结果]
收盘价相对误差超过 $0.1\%$ 的比例为 \textbf{0.000\%}。
两个完全独立的采集通道给出了一致的价格，这同时验证了双方的数据质量。
\end{keybox}

\subsection{第二项：Point-in-Time}

这是最容易被忽略、也最致命的一项。财务数据必须按\textbf{公告日}对齐，
而不是报告期。检验方法是看净利润这个字段在哪个交易日发生变化。

以平安银行（000001）为例：

\begin{center}
\begin{tabular}{@{}lll@{}}
\toprule
数值变化日 & 归母净利润（累计，元） & 对应的报告 \\
\midrule
2020-02-14 & $2.8195\times10^{10}$ & 2019 年年报 \\
2020-04-21 & $8.5480\times10^{9}$ & 2020 年一季报 \\
2020-08-28 & $1.3678\times10^{10}$ & 2020 年中报 \\
2020-10-22 & $2.2398\times10^{10}$ & 2020 年三季报 \\
\bottomrule
\end{tabular}
\end{center}

四个日期与平安银行真实公告日完全吻合。\textbf{这份快照的财务数据是 PIT 的}，
可以直接用于建模。浦发银行（600000）的检验结果同样吻合。

\begin{warnbox}[如果这一项没通过]
即便价格、成交量全部正确，用报告期对齐的财务数据也会让模型提前 1--4 个月
"知道"财报内容。这种前视偏差不会报错，只会让回测好看得离谱。
所以接入任何新数据源，这一项必须第一个做。
\end{warnbox}

\subsection{第三项：覆盖偏差}

前两项都通过之后，第三项暴露了问题。

\begin{center}
\begin{tabular}{@{}lrrr@{}}
\toprule
交易日 & 外部快照 & 本地库 & 缺口 \\
\midrule
2022-06-30 & 3,104 & 4,556 & $-1,452$ \\
2024-06-28 & 3,072 & 5,025 & $-1,953$ \\
2026-08-31 & 4,186 & 5,202 & $-1,016$ \\
\bottomrule
\end{tabular}
\end{center}

外部快照的股票池只有全市场的 \textbf{76\%--83\%}，而且\textbf{各年比例不同}
（2021 年日均只有 2,818 只，是六年里最低的）。缺失的既有 ST 股，也有万科A
这样的蓝筹。

\begin{keybox}[这份快照的正确用法]
它是一个\textbf{已过滤的股票池}，不是全市场。因此：
\begin{enumerate}
  \item \textbf{不能}用它替代本地库做全市场横截面研究——股票池的边界随年份变化，
        会直接污染因子排序；
  \item \textbf{可以}用它做三件事：补充真实成交额与换手率、引入资金流因子族、
        提供日频估值；
  \item \textbf{必须}在任何相关结论旁注明覆盖范围。
\end{enumerate}
\end{keybox}

\subsection{第四项：本地估算到底有多不准}

第三项说明了不能替换股票池。但还有一个更基本的问题：这份快照宣称的价值之一是
"提供真实成交额与换手率，而本地只能估算"。那本地到底估得多差？

用 554 万行匹配数据实测：

\begin{center}
\begin{tabular}{@{}llrrr@{}}
\toprule
字段 & 本地的做法 & 中位误差 & P90 & P99 \\
\midrule
成交额 & 成交量 $\times$ 当日均价 & $0.31\%$ & $1.03\%$ & $2.37\%$ \\
流通股本 & 月度快照前向填充 & $0.000\%$ & $0.022\%$ & $21.9\%$ \\
换手率 & 成交量 $/$ 流通股本 & $\approx 1\%$ & -- & -- \\
\bottomrule
\end{tabular}
\end{center}

成交额的估算中位误差 $0.31\%$、99 分位 $2.37\%$，这个精度对流动性因子
（Amihud 非流动性、成交额对数、流动性门槛）完全够用。
流通股本在 96\% 的交易日上与真值一致，只有 1.1\% 的行偏差超过 20\%，
集中在送转当年。

\begin{keybox}[这笔交易不划算]
"用真实成交额替换估算值"这个卖点，在本项目里收益接近零（$0.3\%$ 的精度提升），
而代价是引入 20\% 的样本缺口与一个逐年变化的股票池边界。
\textbf{为了 $0.3\%$ 的精度牺牲三成小盘股，不划算。}

这也是一个通用教训：接入新数据源时，先量化"我原来有多差"，
再决定要不要替换。很多人换数据源换掉的不是误差，而是可复现性。
\end{keybox}

\section{资金流因子族：一次"换皮检验"的完整演示}

这是本章最有教学价值的部分。资金流数据看起来信息量很大，
但必须先回答一个问题：\textbf{它是不是"当日涨跌幅"或"换手率"的换皮？}

A 股的主力净流入与当日涨跌高度同向，而涨跌本身已经是反转因子的原料。
如果不做检验就直接把 33 个字段丢进模型，得到的很可能是"更贵的反转因子"。

\subsection{因子构造}

所有资金流都除以当日成交额，化为无量纲的"净流入强度"：

\begin{align}
  \text{mfr} &= \frac{\text{主力净流入}}{\text{成交额}}, \qquad
  \text{sfr} = \frac{\text{超大单净流入}}{\text{成交额}}, \\
  \text{lfr} &= -\frac{\text{小单净流入}}{\text{成交额}}, \\
  \text{mfr\_20} &= \frac{\sum_{20}\text{主力净流入}}{\sum_{20}\text{成交额}}, \\
  \text{mfr\_pos\_20} &= \frac{1}{20}\sum_{20}\mathbb{1}[\text{mfr} > 0].
\end{align}

\subsection{统一标签与统计口径后的复核}
旧版此处误用收盘到收盘标签。现在与主面板共用次日复权开盘至第六日复权开盘收益，
即 $y_t=O^{adj}_{t+6}/O^{adj}_{t+1}-1$。标签按个股观测行位移，缺失交易日仍需另做日历对齐审计。
调整后 IC 控制五日反转与二十日换手；没有把这些新因子加入主模型重训。
{\small
\begin{longtable}{@{}lrrrrrr@{}}
\toprule
因子 & 原始 IC & 调整后 IC & 原始 $t$ & 调整后 $t$ & $n$ & $n_{\rm eff}$ \\
\midrule
\endhead
mfr\_pos\_20 & -0.0151 & 0.0246 & -4.68 & 5.41 & 1598 & 386 \\
mfr\_chg & -0.0060 & 0.0006 & -2.09 & 0.29 & 1598 & 522 \\
sfr & -0.0061 & 0.0004 & -3.88 & 0.28 & 1609 & 1187 \\
bfr & -0.0045 & -0.0029 & -2.40 & -1.89 & 1609 & 962 \\
mfr\_stab\_20 & -0.0164 & -0.0053 & -4.52 & -1.66 & 1598 & 400 \\
mfr & -0.0108 & -0.0061 & -5.26 & -3.71 & 1609 & 1097 \\
mfr\_20 & -0.0194 & -0.0080 & -5.70 & -2.75 & 1598 & 408 \\
lfr & -0.0117 & -0.0083 & -5.20 & -4.29 & 1609 & 973 \\
\bottomrule
\end{longtable}
}

原始 IC 与调整后 IC 均由完整日频 IC 序列汇总；$h=5$ 的日频标签重叠，
均值检验使用 Bartlett Newey--West 协方差，滞后阶数为
$L=\max(4,\lfloor4(n/100)^{2/9}\rfloor)$。$n$ 是有值的 IC 日期数，
$n_{\rm eff}=n\gamma_0/\widehat\Omega$ 是原始 IC 均值方差等价的样本量，
不是独立财报次数，也不是未经假设可确定的真实样本数。
完整表另含调整后有效样本量；不同带宽及每五日采样的敏感性见配套 CSV。
本表不再平均各年的 $t$ 值。

\subsection{怎么读这张表}
\texttt{mfr} 的原始 IC 为 \AuditFlowMfr{}；\texttt{mfr\_pos\_20} 的调整后 IC 为
\AuditFlowPos{}，HAC $t=\AuditFlowPosT$。这些数字替换旧收盘标签的结论。
不能仅凭与反转相关就认定某因子“只有换皮”，也不能仅凭残差 IC 显著就证明独立可交易收益。
残差化控制变量、股票池、成本、多重筛选和时间稳定性都会影响解释。
本轮对八个因子全部保留结果，没有事后只展示最优一个。

\section{本章小结}
补充库覆盖 2020--2026 的变化股票池，不能替代主库。
统一开盘标签提高了定义的一致性，但股票池、区间和控制变量仍不同，跨库结果不能直接排名。
价格匹配与个别公告日核验是局部证据；2021 年整年缺失证券及筛选规则仍待解释。
